% =====================================================================
%  P1 small submission version.
%  Compiles standalone with pdflatex; run twice for references.
% =====================================================================
\documentclass[11pt,a4paper]{article}

\usepackage{amsmath}
\usepackage{amssymb}
\usepackage{amsthm}
\usepackage[margin=1in]{geometry}
\usepackage[colorlinks=true,linkcolor=blue,citecolor=blue]{hyperref}

\theoremstyle{plain}
\newtheorem{theorem}{Theorem}
\newtheorem{corollary}[theorem]{Corollary}
\newtheorem{proposition}[theorem]{Proposition}
\newtheorem{lemma}[theorem]{Lemma}
\newtheorem{conjecture}[theorem]{Conjecture}
\theoremstyle{definition}
\newtheorem{observation}[theorem]{Observation}
\newtheorem{measurement}[theorem]{Measurement}
\theoremstyle{remark}
\newtheorem{note}[theorem]{Note}

\renewcommand{\SS}{\mathfrak{S}}
\newcommand{\AAA}{\mathfrak{A}}
\newcommand{\Emu}{E_\mu}
% T_m is the endpoint of the inner sum, not the weighted sum T_w;
% at m = 1 the two symbols would otherwise collide.
\newcommand{\Tcut}{\mathcal T}
\DeclareMathOperator{\rad}{rad}
\DeclareMathOperator{\lcm}{lcm}

% typographic slack: let TeX stretch a line rather than run into the margin
\emergencystretch=3em

\title{A fixed-class M\"obius-weighted estimate and a residual relation in the Huang--Li Goldbach reduction}
\author{Youngrok Lee\\
\small Independent researcher, Republic of Korea}
\date{}

\begin{document}
\maketitle

\begin{abstract}
Huang and Li \cite{HL} obtain a conditional binary Goldbach conclusion for
sufficiently large even integers from the Elliott--Halberstam conjecture
together with a M\"obius-twisted variant.  In their Section~3 the
M\"obius-twisted hypothesis is used to control two signed fixed-class terms,
$E_4(\alpha)$ and $E_3(\alpha)$.  In the regime corresponding to their
corollary, we prove that the fixed-class sum underlying $E_4$ is
$\ll_A N(\log N)^{-A}$ for every $A>0$, uniformly in the truncation
point, using divisor switching and Bombieri--Vinogradov.  Thus that use of
the M\"obius-twisted hypothesis can be replaced by an unconditional
estimate.  The proof completes the fixed divisor sum and then works on the
complementary range, where the surviving M\"obius factor lies on a
cofactor within the Bombieri--Vinogradov range.  This mechanism concerns
the signed fixed-class combination only and does not establish the full
$EH_\mu$ hypothesis.  For the remaining logarithmically weighted term we derive
\[
E_3(\alpha)=\widetilde r(N)-\mathfrak S(N)
\Bigl(N-\sum_{n<N}\Lambda(n)\mu(N-n)\Bigr)
+O_A\!\bigl(N(\log N)^{-A}\bigr).
\]
Consequently, at this error scale, the $E_3$ estimate used in the
Huang--Li decomposition considered here is equivalent to their
equation~(22); this is a relation inside that reduction, not a new
estimate for the remaining correlation.  We also record the moving-cut correction to their published
equation~(18), following the prior defect report and Moya's repair
\cite{Moya}, and bound the corresponding outside correction term.  No
unconditional progress toward Goldbach is claimed.
\end{abstract}


% =====================================================================
\section{Introduction}
\label{sec:intro}

Let $N$ be a large even integer, $\Lambda$ the von Mangoldt function and
$\mu$ the M\"obius function.  Put
\[
  \widetilde r(N)=\sum_{n<N}\Lambda(n)\Lambda(N-n),\qquad
  \SS(N)=2\prod_{p>2}\Bigl(1-\frac1{(p-1)^2}\Bigr)
  \prod_{\substack{p\mid N\\p>2}}\Bigl(1+\frac1{p-2}\Bigr).
\]
Write $EH_\mu(Q)$ for the assertion that, for every $A>0$,
\begin{equation}\label{eq:EHmu}
  \sum_{q\le Q}\max_{y<N}\max_{(a,q)=1}
  \left|\sum_{\substack{n\le y\\n\equiv a\ (q)}}
  \Lambda(n)\mu(N-n)-\frac1{\varphi(q)}
  \sum_{n\le y}\Lambda(n)\mu(N-n)\right|
  \ll_A \frac{N}{(\log N)^A}.
\end{equation}
Huang and Li \cite{HL} state, under their Elliott--Halberstam
hypotheses, the lower bound
\begin{equation}\label{eq:HLthm1}
  \widetilde r(N)\ge \SS(N)(1-\AAA(N))N
  +O\!\left(\frac{N}{(\log N)^A}\right),
\end{equation}
where
\begin{equation}\label{eq:AN}
  \AAA(N)=\prod_{\substack{q\nmid N\\q>2}}
  \left(1-\frac1{q(q-1)}\right).
\end{equation}
Their corollary uses Bombieri--Vinogradov for the ordinary prime
progression input and leaves $EH_\mu(N^{\theta'})$ with one fixed
$\theta'>1/2$.

The M\"obius-twisted hypothesis is used in their Section~3 to bound two
fixed-class error terms.  With $K=(N-1)/\alpha$ they are
\begin{align}
 E_3(\alpha)&=\sum_{\substack{k<K\\(k,N)=1}}\mu(k)\log k
 \left[\sum_{\substack{n<N\\n\equiv N\ (k)}}\Lambda(n)\mu(N-n)
 -\frac1{\varphi(k)}\sum_{n<N}\Lambda(n)\mu(N-n)\right],
 \label{eq:E3}\\
 E_4(\alpha)&=\sum_{\substack{k<K\\(k,N)=1}}\mu(k)
 \Biggl[\sum_{\substack{n<N\\n\equiv N\ (k)}}
 \Lambda(n)\mu(N-n)\log(N-n)\notag\\
 &\hspace{8em}-\frac1{\varphi(k)}\sum_{n<N}
 \Lambda(n)\mu(N-n)\log(N-n)\Biggr].
 \label{eq:E4}
\end{align}
For $(k,N)=1$ and $1\le t<N$ define
\begin{equation}\label{eq:Emu}
 \Emu(t;k)=\sum_{\substack{n\le t\\n\equiv N\ (k)}}
 \Lambda(n)\mu(N-n)-\frac1{\varphi(k)}\sum_{n\le t}
 \Lambda(n)\mu(N-n),\qquad
 C(t)=\sum_{n\le t}\Lambda(n)\mu(N-n),
\end{equation}
and
\[
 T_w(t)=\sum_{\substack{k<K\\(k,N)=1}}\mu(k)w(k)\Emu(t;k).
\]
For endpoint notation set $C(N):=C(N-1)$ and
$\Emu(N;k):=\Emu(N-1;k)$. Thus $E_3=T_{\log}(N-1)$, while $E_4$ is obtained from $T_1(t)$ by
Abel summation.  We use the \emph{Corollary-1 regime}
$\alpha=N^{1-\theta'}$, $\theta'\in(1/2,1)$.  Then
$K=(N-1)N^{\theta'-1}=N^{\theta'}+O(1)$ and
$N/K\asymp N^{1-\theta'}$.

\begin{theorem}\label{thm:A}
Fix $\theta'\in(1/2,1)$.  Uniformly for real $K$ satisfying
$|K-N^{\theta'}|\le1$ and every $A>0$,
\[
 \sup_{1\le t<N}
 \left|\sum_{\substack{k<K\\(k,N)=1}}\mu(k)\Emu(t;k)\right|
 \ll_{A,\theta'}\frac{N}{(\log N)^A}.
\]
All terms in the proof other than the Bombieri--Vinogradov contribution
are $O(Ne^{-c\sqrt{\log N}})$.
\end{theorem}

The theorem concerns a signed sum in the fixed class $a=N$; it is not an
instance of $EH_\mu$, which places absolute values outside each modulus
and maximises over the residue class.

\begin{corollary}\label{cor:B}
In the Corollary-1 regime,
$E_4(\alpha)\ll_A N(\log N)^{-A}$ for every $A>0$ without assuming
$EH_\mu$.  Hence the use of Huang--Li's Lemma~5 to control $E_4$ can be
replaced by Theorem~\ref{thm:A}.  After the moving-cut issue of
Section~\ref{sec:delta} is accounted for, the remaining use of
$EH_\mu$ within the $E_3/E_4$ decomposition considered here is the
estimate for $E_3(\alpha)$.
\end{corollary}

\begin{theorem}\label{thm:C}
In the Corollary-1 regime, for every $A>0$,
\begin{equation}\label{eq:Crelation}
 E_3(\alpha)=\widetilde r(N)-\SS(N)\bigl(N-C(N)\bigr)
 +O_A\!\left(\frac{N}{(\log N)^A}\right).
\end{equation}
\end{theorem}

\begin{corollary}[conditional consequences]\label{cor:Cconseq}
Assume in addition that
\[
 E_3(\alpha)\ll_A N(\log N)^{-A}
\]
for every $A>0$.  This is the remaining input in the Huang--Li
$E_3/E_4$ decomposition and is not proved here.  Then
\eqref{eq:Crelation} is, at the same error scale, equivalent to their
equation~(22).  Together with the unconditional Bombieri--Vinogradov
estimate used by Huang--Li,
\[
 |C(N)|\le U(N),\qquad
 U(N)=\AAA(N)N+O_B\!\left(N(\log N)^{-B}\right),
\]
it yields binary Goldbach for all sufficiently large even $N$.
Under the same additional estimate,
$\widetilde r(N)\sim\SS(N)N$ if and only if $C(N)=o(N)$.
\end{corollary}

Theorem~\ref{thm:C} is specific to the logarithmic weight occurring in
this reduction: completing the divisor sum invokes
$\mu*\log=\Lambda$ and reproduces $\widetilde r(N)$.  Corollary~\ref{cor:Cconseq}
is only a conditional consequence of the additional $E_3$ estimate;
that estimate is not supplied by Theorem~\ref{thm:C}.  We do not claim
that every alternative smoothing or reformulation is covered by this
calculation.

\begin{proposition}[one-sided sufficient condition]\label{prop:onesided}
In the Corollary-1 regime, binary Goldbach for all sufficiently large even
$N$ follows if, for one fixed $\varepsilon>0$,
\begin{equation}\label{eq:onesided}
 E_3(\alpha)>-(1-\varepsilon)\SS(N)(1-\AAA(N))N.
\end{equation}
\end{proposition}

\begin{proposition}[absolute-value sufficient condition]\label{prop:nolog}
Put
\begin{equation}\label{eq:Bsum}
 B_\mu(N)=\sum_{\substack{k<K\\(k,N)=1}}
 \mu^2(k)(\log k)|\Emu(N;k)|.
\end{equation}
Then binary Goldbach for all sufficiently large even $N$ follows if, for
one fixed $\varepsilon>0$,
\begin{equation}\label{eq:nolog}
 B_\mu(N)\le(1-\varepsilon)\SS(N)(1-\AAA(N))N.
\end{equation}
The full hypothesis $EH_\mu(N^{\theta'})$ implies this condition, but
the converse is not asserted.
\end{proposition}

The scope is deliberately limited.  Theorem~\ref{thm:A} gives no new
bound for $C(N)$, and Theorem~\ref{thm:C} converts the surviving
$E_3$ estimate into a relation containing the Goldbach sum rather than
proving that estimate.  Accordingly no unconditional progress toward the
Goldbach conjecture is claimed.  Section~\ref{sec:delta} concerns a
published rearrangement in \cite{HL}; it provides a correction and does
not claim that their stated conditional theorem is false.

% =====================================================================
\section{Notation and the mechanism}\label{sec:mechanism}

$p$ and $q$ always denote primes. $\varphi$ is Euler's function,
$\tau_3$ the $3$-fold divisor function, $\rad(u)=\prod_{p\mid u}p$.
Constants $c,c_1,\dots$ are positive and absolute unless subscripted;
implied constants may depend on $A$ and $\theta'$. The letter $A$ is
reserved for the free exponent of \eqref{eq:EHmu}; the local factor of
\eqref{eq:AN} is $\AAA$.

The mechanism of Theorem~\ref{thm:A} is the following. The residue
class in \eqref{eq:Emu} is the fixed class $n\equiv N\pmod k$, i.e.
\[
  n\equiv N\pmod{k}\quad\Longleftrightarrow\quad k\mid N-n .
\]
So the $k$-sum is a divisor sum over $u=N-n$, namely the incomplete
sum $\sigma_K(u)=\sum_{k\mid u,\ k<K,\ (k,N)=1}\mu(k)$.  Put
$\delta=(\theta'-\tfrac12)/2>0$; then
$N/K\ll N^{1/2-2\delta}$.

\emph{First} one \textbf{completes} it. By Lemma~\ref{lem:complete},
for squarefree $u$ the complete sum
$\sum_{k\mid u,\,(k,N)=1}\mu(k)$ equals $\mathbf 1_{\rad(u)\mid N}$,
which forces $u\mid\rad(N)$ and so leaves only $N^{o(1)}$ terms.
\emph{Then} the work sits entirely in the complementary sum, over
$k\ge K$ --- and it is only on that range that writing $u=mk$ gives
\[
  m\;=\;u/k\;<\;N/K\;\ll\;N^{1/2-2\delta}.
\]
\emph{Finally}, for squarefree $u=mk$ the factorisation forces
$(m,k)=1$ and $\mu(u)\mu(k)=\mu(m)\mu^2(k)$: the M\"obius factor on
the \emph{long} variable $k$ cancels against itself and leaves the
nonnegative $\mu^2$, while the surviving M\"obius sits on the
\emph{short} variable $m$. That is the classical ``M\"obius on the
short variable plus Bombieri--Vinogradov'' configuration.

The completion step is essential for this proof.  Before completion the
cofactor $u/k$ need not be $O(N^{1-\theta'})$, so the surviving
M\"obius factor is not placed on the range controlled by the
Bombieri--Vinogradov step below.

For the present theorem, the only size fact needed here is
$m\ll N^{1/2-2\delta}$, which places the cofactor inside the
Bombieri--Vinogradov range used below. Steps~1--3 of
Section~\ref{sec:proof} carry out the argument in that order.



% =====================================================================
\section{Proof of \texorpdfstring{Theorem~\ref{thm:A}}{the first theorem}}
\label{sec:proof}

Fix $A>0$ and set
\begin{equation}\label{eq:trunc}
  D_0\;=\;E_0\;:=\;\exp\bigl(\sqrt{\log N}\bigr).
\end{equation}
The exponential truncation is independent of $A$ and puts every
non--Bombieri--Vinogradov tail below every power of $\log N$, while
remaining harmless for the modulus bound \eqref{eq:level}.

\subsection{Step 0: the subtracted mean term}

By \eqref{eq:Emu},
\begin{equation}\label{eq:split0}
  T_1(t)\;=\;D(t)\;-\;C(t)\,B(K),
  \qquad
  D(t):=\sum_{\substack{k<K\\(k,N)=1}}\mu(k)
        \sum_{\substack{n\le t\\ n\equiv N\ (k)}}\Lambda(n)\mu(N-n),
\end{equation}
where $B(K)=\sum_{k<K,(k,N)=1}\mu(k)/\varphi(k)$. Huang--Li's Lemma~1
(a form of Goldston--Y{\i}ld{\i}r{\i}m \cite{GY}) gives the same
bound with $k\le K$; if $K$ is an integer the endpoint difference is
$O(K^{-1+o(1)})$ and is absorbed. Hence
$B(K)\ll e^{-c_1\sqrt{\log K}}$, and trivially
$|C(t)|\le\sum_{n<N}\Lambda(n)\ll N$. Hence
\begin{equation}\label{eq:meanterm}
  |C(t)B(K)|\;\ll\;Ne^{-c\sqrt{\log N}}
\end{equation}
uniformly in $t$.

\subsection{Step 1: divisor switching (an exact identity)}

Since $n\equiv N\pmod k$ with $n\le t<N$ is the same as $u:=N-n$ being
a multiple of $k$ in $[N-t,N)$,
\begin{equation}\label{eq:switch}
  D(t)\;=\;\sum_{N-t\le u<N}\Lambda(N-u)\,\mu(u)\,\sigma_K(u),
  \qquad
  \sigma_K(u):=\sum_{\substack{k\mid u,\ k<K\\ (k,N)=1}}\mu(k).
\end{equation}
This is a finite rearrangement and hence exact.

\subsection{Step 2: the complete divisor sum}

\begin{lemma}\label{lem:complete}
For every $u\ge1$,
$\displaystyle\sum_{k\mid u,\ (k,N)=1}\mu(k)=\mathbf 1_{\rad(u)\mid N}$.
\end{lemma}

\begin{proof}
The sum is multiplicative in $u$ with local factor $1$ at $p\mid N$
and $1-1=0$ at $p\nmid N$. Hence it vanishes unless every prime of $u$
divides $N$.
\end{proof}

Splitting $\sigma_K(u)$ into the complete sum minus the tail $k\ge K$,
\eqref{eq:switch} becomes $D(t)=P(t)-R(t)$ with
\begin{equation}\label{eq:PR}
  P(t)=\!\!\sum_{\substack{N-t\le u<N\\ \rad(u)\mid N}}\!\!\Lambda(N-u)\mu(u),
  \qquad
  R(t)=\!\!\sum_{N-t\le u<N}\!\!\Lambda(N-u)\mu(u)
       \sum_{\substack{k\mid u,\ k\ge K\\ (k,N)=1}}\mu(k).
\end{equation}
In $P(t)$ the conditions $\mu(u)\ne0$ and $\rad(u)\mid N$ force
$u\mid\rad(N)$, so there are at most $2^{\omega(N)}=N^{o(1)}$ terms,
each $\ll\log N$:
\begin{equation}\label{eq:P}
  P(t)\ll N^{o(1)}.
\end{equation}

\subsection{Step 3: the residual, and the degeneracy lemma}

Write $u=mk$ with $k\ge K$, so $m=u/k<N/K=:M$. For squarefree $u$ the
factorisation forces $(m,k)=1$ and $\mu(u)\mu(k)=\mu(m)\mu^2(k)$, so
\begin{equation}\label{eq:R1}
  R(t)=\sum_{m<M}\mu(m)
      \sum_{\substack{k\ge K,\ (k,m)=1,\ (k,N)=1\\ N-t\le mk<N}}
      \mu^2(k)\,\Lambda(N-mk).
\end{equation}

\begin{lemma}[degeneracy]\label{lem:degen}
The contribution to \eqref{eq:R1} of the terms with $(k,N)>1$ is
$O(N^{o(1)})$. The same bound holds for the terms in which the modulus
constructed in Step~4 is not coprime to $N$.
\end{lemma}

\begin{proof}
Suppose $p\mid(k,N)$ and $\Lambda(N-mk)\ne0$, say $N-mk=q^\ell$. Since
$p\mid k$ and $p\mid N$ we get $p\mid N-mk$, hence $q=p$ and
$n:=N-mk=p^\ell$ with $p\mid N$. The number of such $n<N$ is
$\le\sum_{p\mid N}\log N/\log p\ll(\log N)^2$; for each, the pairs
$(m,k)$ with $mk=N-n$ number $\le\tau(N-n)\ll N^{o(1)}$, and each term
is $\ll\log N$.

For the second sentence, let $q=m\lcm(d^2,e)$ be the modulus built in
Step~4 and suppose $g=(q,N)>1$. A term of \eqref{eq:R2} in that class
has $n\equiv N\pmod q$ with $1\le n\le\Tcut_m(t)$, so $g\mid N-n$;
together with $g\mid N$ this gives $g\mid n$, and any $p\mid g$ then
divides $n$, so $\Lambda(n)\ne0$ forces $n=p^\ell$ with $p\mid N$. The
number of such $n<N$ is again $\le\sum_{p\mid N}\log N/\log p
\ll(\log N)^2$. For each such $n$, a triple $(m,d,e)$ places it in a
degenerate class only if $m\lcm(d^2,e)\mid N-n$, so each of $m$, $d^2$
and $e$ divides $N-n$ and the triples number
$\ll\tau(N-n)^3\ll N^{o(1)}$; with $\Lambda(n)\ll\log N$ the total is
$\ll N^{o(1)}$.
\end{proof}

By Lemma~\ref{lem:degen} we may drop the condition $(k,N)=1$ from
\eqref{eq:R1} at a cost of $O(N^{o(1)})$. This is essential: expanding
$\mathbf 1_{(k,N)=1}$ by M\"obius over $e\mid N$ would multiply the
number of Bombieri--Vinogradov calls by $2^{\omega(N)}=N^{o(1)}$ and
destroy the budget.

\subsection{Step 4: unfolding \texorpdfstring{$\mu^2$}{mu-squared} and the
coprimality}

Insert $\mu^2(k)=\sum_{d^2\mid k}\mu(d)$ and
$\mathbf 1_{(k,m)=1}=\sum_{e\mid(k,m)}\mu(e)$ and truncate at
$D_0,E_0$ of \eqref{eq:trunc}.

Two indices are truncated at once, so the discarded set has to be named
before it is bounded, and it is not the union of two independent tails.
Write the double sum over $(d,e)$ and split it as
\begin{equation}\label{eq:twotails}
  \sum_{d,e}\;=\;\sum_{d\le D_0}\sum_{e\le E_0}
  \;+\;\underbrace{\sum_{d>D_0}\sum_{e}}_{\text{$e$-sum complete}}
  \;+\;\underbrace{\sum_{d\le D_0}\sum_{e>E_0}}_{\text{$d$-sum truncated}},
\end{equation}
an exact identity. In the second block the inner $e$-sum is complete,
so it collapses to $\mathbf 1_{(k,m)=1}$. In the third block the
$d$-sum is truncated and must be bounded as such; no cancellation from
the complete identity for $\mu^2(k)$ is used below.

\emph{Tail $d>D_0$.} Such terms have $d^2\mid k$ and $mk<N$, so
necessarily $d\le\sqrt{k}<\sqrt{N/m}$.  Also
$n=N-mk\equiv N\pmod{md^2}$; the number of $n\le N$ in that
progression is $\ll N/(md^2)+1$ and each $\Lambda\ll\log N$. Hence
\[
  \ll\log N\sum_{m<M}\sum_{D_0<d\le\sqrt{N/m}}
       \Bigl(\frac{N}{md^2}+1\Bigr)
  \ll\frac{N(\log N)^2}{D_0}+\sqrt{NM}\log N,
\]
which is
\[
  \ll N(\log N)^2e^{-\sqrt{\log N}}+N^{1-\theta'/2}\log N.
\]

\emph{Tail $e>E_0$.} Here $m$ is squarefree, hence every $e\mid m$
is squarefree.  Taking absolute values before summing over the truncated
$d$-range gives moduli $m\lcm(d^2,e)$.  Since
\[
 \frac1{\lcm(d^2,e)}=\frac{(d,e)}{d^2e}
 \quad(e\ \hbox{squarefree}),\qquad
 \sum_{d\ge1}\frac{(d,e)}{d^2}\ll\prod_{p\mid e}\Bigl(1+\frac1p\Bigr)
 \ll\log\log(3e),
\]
the progression-counting part is
\[
 \ll N\log N\sum_{m<M}\frac1m
       \sum_{\substack{e\mid m\\e>E_0}}
       \frac{\log\log(3e)}e
 \ll \frac{N(\log N)^3\log\log N}{E_0}.
\]
The ``$+1$'' in the progression count contributes at most
\[
 \ll D_0\log N\sum_{m<M}\tau(m)
 \ll D_0M(\log N)^2.
\]
Thus this tail is
$\ll N(\log N)^3\log\log N\,e^{-\sqrt{\log N}}
   +M e^{\sqrt{\log N}}(\log N)^2$.

Both are $\ll Ne^{-c\sqrt{\log N}}$ for any $c<1$, hence below every
power of $\log N$. In the remaining range put
$L:=\lcm(d^2,e)$ and $q:=mL$; the conditions $d^2\mid k$, $e\mid k$
and $n=N-mk$ give $n\equiv N\pmod q$, and $mk<N$,
$mk\ge\max(N-t,mK)$ give $1\le n\le\Tcut_m(t)$ where
\begin{equation}\label{eq:Tm}
  \Tcut_m(t):=\min\bigl(t,\ N-mK\bigr).
\end{equation}
Therefore, with $\psi(y;q,a)=\sum_{n\le y,\,n\equiv a\,(q)}\Lambda(n)$,
\begin{equation}\label{eq:R2}
  R(t)=\sum_{m<M}\mu(m)\sum_{d\le D_0}\mu(d)
         \sum_{\substack{e\mid m\\ e\le E_0}}\mu(e)\,
         \psi\bigl(\Tcut_m(t);\,q,\,N\bigr)
         \;+\;O\!\left(Ne^{-c\sqrt{\log N}}\right).
\end{equation}
The error here is the two truncation tails just bounded, and it is
exponential: Bombieri--Vinogradov has not been applied yet, and the
second sentence of Theorem~\ref{thm:A} --- that every ingredient but
Bombieri--Vinogradov saves exponentially --- is read off this line
among others.
By Lemma~\ref{lem:degen} we may further restrict to $(q,N)=1$ at a cost
of $O(N^{o(1)})$; the lemma obtains this by counting the admissible $n$,
of which there are $\ll(\log N)^2$, rather than the parameter triples.



Note the size of the modulus: $q=m\lcm(d^2,e)\le MD_0^2E_0$, so by
\eqref{eq:trunc}
\begin{equation}\label{eq:level}
  q\;\le\;M e^{3\sqrt{\log N}}
  \;\ll\;N^{1/2-2\delta}e^{3\sqrt{\log N}}
  \;\le\;N^{1/2-\delta}
\end{equation}
for $N$ large. This is the level at which Bombieri--Vinogradov is
applied.

\subsection{Step 5: Bombieri--Vinogradov}

\begin{lemma}[$\tau$-weighted Bombieri--Vinogradov]\label{lem:BV}
For all $A,B>0$ there is $C=C(A,B)$ such that, for
$Q\le N^{1/2}(\log N)^{-C}$,
\[
  \sum_{q\le Q}\tau_3(q)^{B}\ \max_{y\le N}\ \max_{(a,q)=1}
  \Bigl|\psi(y;q,a)-\frac{y}{\varphi(q)}\Bigr|
  \;\ll_{A,B}\;\frac{N}{(\log N)^{A}}.
\]
\end{lemma}

\begin{proof}
This is a standard weighted consequence of the maximal
Bombieri--Vinogradov theorem; see, for example,
\cite[\S20.2]{MoV2}.  Put
\[
 \mathcal E_q:=\max_{y\le N}\max_{(a,q)=1}
 \Bigl|\psi(y;q,a)-\frac{y}{\varphi(q)}\Bigr|.
\]
By Cauchy--Schwarz,
\[
 \sum_{q\le Q}\tau_3(q)^B\mathcal E_q
 \le
 \Bigl(\sum_{q\le Q}\tau_3(q)^{2B}\mathcal E_q\Bigr)^{1/2}
 \Bigl(\sum_{q\le Q}\mathcal E_q\Bigr)^{1/2}.
\]
The trivial estimate
$\mathcal E_q\ll (N/\varphi(q))\log N$ and the standard mean-value
bounds for fixed powers of the divisor function give
\[
 \sum_{q\le Q}\tau_3(q)^{2B}\mathcal E_q
 \ll_B N(\log N)^{C_1(B)}.
\]
On the other hand, the maximal Bombieri--Vinogradov theorem gives, for
every $D>0$, a constant $C_2(D)$ such that
\[
 \sum_{q\le Q}\mathcal E_q\ll_D \frac{N}{(\log N)^D}
 \qquad
 \bigl(Q\le N^{1/2}(\log N)^{-C_2(D)}\bigr).
\]
Taking $D=2A+C_1(B)+2$ and increasing the modulus exponent if necessary
proves the stated estimate.  No strengthened distribution theorem is
being used here; only the usual Bombieri--Vinogradov theorem with an
arbitrarily large logarithmic saving.
\end{proof}



A modulus $q$ arises in \eqref{eq:R2} from at most
$\sum_{m\mid q}\tau(q/m)^2\le\tau(q)^3\le\tau_3(q)^3$ triples
$(m,d,e)$: given $m\mid q$ there are at most $\tau(q/m)$ choices of
$d$ with $d^2\mid q/m$ and at most $\tau(q/m)$ of $e$ with
$\lcm(d^2,e)=q/m$. Writing $\psi(\Tcut_m;q,N)=\Tcut_m/\varphi(q)+\mathcal E$
and applying Lemma~\ref{lem:BV} with $B=3$ and \eqref{eq:level},
\begin{equation}\label{eq:R3}
  R(t)\;=\;\mathrm{MT}(t)\;+\;O\!\left(\frac{N}{(\log N)^{A}}\right),
  \qquad
  \mathrm{MT}(t)=\sum_{\substack{m<M\\ (m,N)=1}}
       \mu(m)\,\Tcut_m(t)\,c_{D_0,E_0}(m),
\end{equation}
where
$c_{D_0,E_0}(m)=\sum_{d\le D_0}\sum_{e\mid m,\,e\le E_0}
\mu(d)\mu(e)\mathbf 1_{(d,N)=1}/\varphi(m\lcm(d^2,e))$.

\subsection{Step 6: the density}

\begin{lemma}\label{lem:density}
Let $m$ be squarefree with $(m,N)=1$ and put
\[
  c(m):=\sum_{d\ge1}\ \sum_{e\mid m}
     \frac{\mu(d)\mu(e)\mathbf 1_{(d,N)=1}}{\varphi(m\lcm(d^2,e))}.
\]
Then $c(m)=\AAA(N)\lambda(m)/m$ with $\AAA(N)$ as in \eqref{eq:AN} and
$\lambda(m)=\prod_{p\mid m}\bigl(1-\tfrac{1}{p(p-1)}\bigr)^{-1}$.
Moreover the two truncations contribute differently:
\begin{equation}\label{eq:ctrunc}
  c_{D_0,E_0}(m)\;=\;c(m)
  \;+\;O\!\left(\frac{d(m)}{\varphi(m)D_0}\right)
  \;+\;O\!\left(\frac{d(m)\log\log M}{\varphi(m)E_0}\right).
\end{equation}
The divisor factor is retained in both truncation errors; it is harmless
only after the subsequent sum over $m$.
\end{lemma}

\begin{proof}
The sum is multiplicative. If $p\nmid m$ then $p\nmid e$ and the local
factor is $1-1/\varphi(p^2)=1-1/(p(p-1))$ for $p\nmid N$, and $1$ for
$p\mid N$. If $p\mid m$ (so $p\nmid N$) the $p$-part of
$m\lcm(d^2,e)$ is $p^{1+\max(2v_p(d),v_p(e))}$, and the four choices
$(v_p(d),v_p(e))\in\{0,1\}^2$ contribute
\[
  \frac{1}{p-1}-\frac{1}{p(p-1)}-\frac{1}{p^2(p-1)}+\frac{1}{p^2(p-1)}
  \;=\;\frac{1}{p}.
\]
Hence $c(m)=\prod_{p\mid m}p^{-1}\cdot
\prod_{p\nmid m,\ p\nmid N}(1-\tfrac{1}{p(p-1)})
=m^{-1}\AAA(N)\lambda(m)$.

The two truncations are done separately.  For $d>D_0$ we retain the
complete $e\mid m$ sum only in absolute value.  Since
$m\lcm(d^2,e)$ is a multiple of $md^2$ and there are at most $d(m)$
choices of $e$,
\[
 \sum_{d>D_0}\sum_{e\mid m}
 \frac1{\varphi(m\lcm(d^2,e))}
 \ll \frac{d(m)}{\varphi(m)}
       \sum_{d>D_0}\frac1{d\varphi(d)}
 \ll \frac{d(m)}{\varphi(m)D_0}.
\]
For $e>E_0$, sum first over $d$.  A prime-by-prime comparison with
$\varphi(me)$ gives
\[
 \sum_{d\ge1}\frac1{\varphi(m\lcm(d^2,e))}
 \ll \frac{\log\log(3e)}{\varphi(me)}
 = \frac{\log\log(3e)}{e\varphi(m)},
\]
since $m$ is squarefree and $e\mid m$, so $\varphi(me)=e\varphi(m)$.  Hence
\begin{equation}\label{eq:etail}
 \sum_{\substack{e\mid m\\e>E_0}}\sum_{d\ge1}
 \frac1{\varphi(m\lcm(d^2,e))}
 \ll \frac{d(m)\log\log M}{\varphi(m)E_0}.
\end{equation}
Finally $\sum_{m<M}d(m)/\varphi(m)\ll(\log M)^2$, while
$D_0^{-1}=E_0^{-1}=e^{-\sqrt{\log N}}$, so both errors are exponentially
small after the $m$-sum in Proposition~\ref{prop:MT}.
\end{proof}

The load-bearing local factor can equivalently be written, for squarefree
$m$, as
\[
  \sum_{g\mid m}\frac{\mu(g)}{\varphi(m/g)\,g\,\varphi(g)}=\frac1m .
\]
Thus the Euler factor is exactly $p^{-1}$ and the $1/\zeta$ factor
appearing in Lemma~\ref{lem:mu} has first order, which is the
cancellation used below.

\begin{lemma}\label{lem:mu}
Let $f(m):=\mu(m)\lambda(m)\mathbf 1_{(m,N)=1}$ and
$G(x):=\sum_{m\le x}f(m)/m$. Then
\[
  G(x)\;\ll\;\frac{N}{\varphi(N)}\,e^{-c\sqrt{\log x}}
   +x^{-1/4}e^{C\sqrt{\log N}}
   \qquad(x\ge2).
\]
\end{lemma}

\begin{proof}
Write $F(s)=\sum_m f(m)m^{-s}=\prod_{p\nmid N}(1-\lambda(p)p^{-s})$
and $F(s)=\zeta(s)^{-1}H(s)$, so $f=\mu*h$ with $h$ multiplicative,
$h(p^j)=1$ for $p\mid N$ and $h(p^j)=1-\lambda(p)=O(p^{-2})$ for
$p\nmid N$. Then
$G(x)=\sum_{b\le x}\frac{h(b)}{b}\sum_{a\le x/b}\frac{\mu(a)}{a}$. For
$b\le\sqrt x$ use $\sum_{a\le y}\mu(a)/a\ll e^{-c\sqrt{\log y}}$, which
is the classical zero-free region estimate \cite{MoV}, together with
$\sum_b|h(b)|/b\ll\prod_{p\mid N}(1-1/p)^{-1}\ll N/\varphi(N)$. For
$b>\sqrt x$ bound
$\sum_{b>y}|h(b)|/b\le y^{-1/2}\sum_b|h(b)|b^{-1/2}
\ll y^{-1/2}\prod_{p\mid N}(1-p^{-1/2})^{-1}\ll y^{-1/2}e^{C\sqrt{\log N}}$.
\end{proof}

Since $M\asymp N^{1-\theta'}$ is a fixed power of $N$, the second term of
Lemma~\ref{lem:mu} is negligible at $x\asymp M$, and
$N/\varphi(N)\ll\log\log N$; so $G(x)\ll e^{-c'\sqrt{\log x}}$ for $x$
of that size. It is not so for small $x$: at $x=2$ the second term is
$2^{-1/4}e^{C\sqrt{\log N}}$, and the bound
$G(x)\ll e^{-c'\sqrt{\log x}}$ fails on the whole of
$x\le e^{4C\sqrt{\log N}}$. Proposition~\ref{prop:MT} integrates $G$
over $[1,M]$ and therefore splits there; the split is carried out in
its proof.

\subsection{Step 7: the main term dies, uniformly in
\texorpdfstring{$t$}{t}}

\begin{proposition}\label{prop:MT}
$\displaystyle\sup_{1\le t<N}|\mathrm{MT}(t)|\ll Ne^{-c\sqrt{\log N}}$.
\end{proposition}

\begin{proof}
By Lemma~\ref{lem:density}, write
\[
\mathrm{MT}(t)=\AAA(N)\sum_{m<M}\frac{f(m)}{m}\Tcut_m(t)
+\mathcal E_{\mathrm{tr}}(t),
\]
with $\Tcut_m(t)$ as in \eqref{eq:Tm}.  The truncation term
$\mathcal E_{\mathrm{tr}}(t)$ is estimated from the two errors in
\eqref{eq:ctrunc} below. Put $m_0:=(N-t)/K$ and $m_*:=\max(1,m_0)$.  On the interval
$[1,M]$, the function $x\mapsto\Tcut_x(t)$ is constant where
$x<m_0$, has derivative $-K$ where $x>m_0$, and satisfies
$\Tcut_M(t)=0$.  Since the summation starts at $m=1$, Abel summation
against $G$ of Lemma~\ref{lem:mu} gives
\[
  \sum_{m<M}\frac{f(m)}{m}\Tcut_m(t)
  = G(M)\,\Tcut_M(t)+K\!\int_{m_*}^{M}\!G(x)\,dx
  = K\!\int_{m_*}^{M}\!G(x)\,dx .
\]
Hence, uniformly in $t$,
\[
  \Bigl|\sum_{m<M}\frac{f(m)}{m}\Tcut_m(t)\Bigr|
  \le K\int_{1}^{M}|G(x)|\,dx .
\]
Lemma~\ref{lem:mu}'s second term is at least $1$ below
$x_0:=e^{4C\sqrt{\log N}}$, so the integral is split there. On
$[1,x_0]$ use the trivial $|G(x)|\ll\log x$, giving
$\ll Kx_0\log x_0\ll N^{\theta'}e^{O(\sqrt{\log N})}$. On $[x_0,M]$ the
two terms of Lemma~\ref{lem:mu} give
\[
  K\!\int_{x_0}^{M}\!\Bigl(\frac{N}{\varphi(N)}e^{-c\sqrt{\log x}}
    +x^{-1/4}e^{C\sqrt{\log N}}\Bigr)dx
  \ll KMe^{-c'\sqrt{\log M}}+KM^{3/4}e^{C\sqrt{\log N}},
\]
and with $KM=N$, $\log M\asymp\log N$ and $M\asymp N^{1-\theta'}$ the two are
$\ll Ne^{-c''\sqrt{\log N}}$ and $\ll N^{(3+\theta')/4}e^{C\sqrt{\log N}}$
respectively. All three pieces are $\ll Ne^{-c\sqrt{\log N}}$ for
$\theta'<1$, the first and the third being powers of $N$ below $N$
itself. The truncation term introduced above has two halves.  Using
\eqref{eq:ctrunc} and $\Tcut_m(t)\le N$ in both gives
\[
 |\mathcal E_{\mathrm{tr}}(t)|
 \ll \frac{N}{D_0}\sum_{m<M}\frac{d(m)}{\varphi(m)}
 +\frac{N\log\log M}{E_0}
   \sum_{m<M}\frac{d(m)}{\varphi(m)}
 \ll N(\log M)^2(1+\log\log M)e^{-\sqrt{\log N}}.
\]
Thus $\mathcal E_{\mathrm{tr}}(t)\ll Ne^{-c\sqrt{\log N}}$ for every fixed $c<1$.
The size here is $N$ and not
$M$: the truncation is per term of the harmonic sum, and the terms are
of the size of $N$.
\end{proof}

The cancellation is genuine and global: the term $m=1$ alone
contributes $\Tcut_1(t)\asymp N$ --- the endpoint of \eqref{eq:Tm}, not
the weighted sum of Theorem~\ref{thm:A} --- so the smallness of
$\mathrm{MT}$ is
entirely due to the cancellation in $\sum f(m)/m$.

\subsection{Conclusion}

Combining \eqref{eq:split0}, \eqref{eq:meanterm}, \eqref{eq:P},
\eqref{eq:R3} and Proposition~\ref{prop:MT}:
\[
  T_1(t)=\underbrace{P(t)}_{N^{o(1)}}
         -\underbrace{\mathrm{MT}(t)}_{\ll Ne^{-c\sqrt{\log N}}}
         -\underbrace{O\!\left(N(\log N)^{-A}\right)}_{\text{BV}}
         -\underbrace{C(t)B(K)}_{\ll Ne^{-c\sqrt{\log N}}},
\]
uniformly in $t<N$. This proves Theorem~\ref{thm:A}.  Every displayed remainder above is
$O(Ne^{-c\sqrt{\log N}})$ except the term supplied by
Lemma~\ref{lem:BV}, which gives the stated arbitrary log-power saving.
\qed

% =====================================================================
\section{Proof of \texorpdfstring{Corollary~\ref{cor:B}}{the corollary}}

Fix $k$ and set
$a_n:=\Lambda(n)\mu(N-n)\bigl(\mathbf 1_{n\equiv N\,(k)}-1/\varphi(k)\bigr)$,
so that $\Emu(t;k)=\sum_{n\le t}a_n$ and the bracket in \eqref{eq:E4}
is $\sum_{n<N}a_n\log(N-n)$. Since $\log(N-n)$ vanishes at $n=N-1$ and
has derivative $-1/(N-t)$, Abel summation gives
\[
  \sum_{n<N}a_n\log(N-n)\;=\;\int_{1}^{N-1}\frac{\Emu(t;k)}{N-t}\,dt .
\]
The weight $\log(N-n)$ does not depend on $k$, so multiplying by
$\mu(k)$ and summing over $k<K$ with $(k,N)=1$ (a finite sum) gives
the exact identity
\[
  E_4(\alpha)\;=\;\int_{1}^{N-1}\frac{T_1(t)}{N-t}\,dt .
\]
Hence $|E_4(\alpha)|\le\bigl(\sup_{t<N}|T_1(t)|\bigr)\log N
\ll_A N(\log N)^{1-A}$ by Theorem~\ref{thm:A}. As $A$ is arbitrary
this is Corollary~\ref{cor:B}.  Thus, in this regime, the invocation
of Huang--Li's Lemma~5 for $E_4$ may be replaced by Theorem~\ref{thm:A};
after the cutoff correction of Section~\ref{sec:delta}, the remaining
use of $EH_\mu$ within the $E_3/E_4$ decomposition considered here is
the estimate for $E_3(\alpha)$. \qed

% =====================================================================
\section{Proof of \texorpdfstring{Theorem~\ref{thm:C}}{the residual relation}}
\label{sec:C}

Take the weight $w_k=\log k$, i.e. the functional $E_3$ of
\eqref{eq:E3}. Steps~0--1 are unchanged, and the complete divisor sum
of Lemma~\ref{lem:complete} is replaced by the following.

\begin{lemma}\label{lem:completelog}
For every $u\ge1$, write $u=u_Nu'$ where $u_N$ is the largest divisor
of $u$ composed of primes dividing $N$. Then
$\sum_{k\mid u,\ (k,N)=1}\mu(k)\log k=-\Lambda(u')$.
\end{lemma}

\begin{proof}
The sum is over $k\mid u'$ and equals $-\Lambda(u')$ by
$\mu*\log=\Lambda$.
\end{proof}

Hence the complete piece of $E_3$ is
\[
  -\sum_{1\le u<N}\Lambda(N-u)\,\mu(u)\,\Lambda(u')
  \;=\;-\sum_{\substack{u<N,\ (u,N)=1}}\Lambda(N-u)\mu(u)\Lambda(u)
        +O(N^{o(1)}\log^2 N),
\]
the discarded terms being those with $(u,N)>1$, which are degenerate in
the sense of Lemma~\ref{lem:degen}: a prime $p\mid(u,N)$ divides $N-u$
as well, so $\Lambda(N-u)\ne0$ forces $N-u$ to be a power of $p$, and
the number of such $u$ is $\ll(\log N)^2$ while each term is
$\ll(\log N)^2$. The set discarded here looks like a set of size $N$
and is not one; that it is not is the same mechanism as
Lemma~\ref{lem:degen}, applied to $u$ instead of to $k$.
Now $\mu(u)\Lambda(u)$ is supported on primes $u=p$, where it equals
$-\log p$; prime powers $u=p^\ell$, $\ell\ge2$, have $\mu(u)=0$. So
the complete piece equals
\begin{equation}\label{eq:goldbachback}
  \sum_{p<N}\Lambda(N-p)\log p\;+\;O\bigl(N^{o(1)}\log^2N\bigr)
  \;=\;\sum_{n<N}\Lambda(n)\Lambda(N-n)
       +O\bigl(N^{1/2}\log^2N\bigr),
\end{equation}
the von--Mangoldt Goldbach convolution, up to the displayed
proper-prime-power error.

Two further terms have to be evaluated, and neither vanishes.

\medskip\noindent
\emph{(i) The residual main term.} Steps~3--6 are repeated with the
extra factor $\log k$, and what has to be checked is that the factor
passes through each of them. It does, for one reason: on the residual
range $k=(N-n)/m$, so that
\begin{equation}\label{eq:logweight}
  \log k\;=\;w_m(n),\qquad w_m(n):=\log\frac{N-n}{m},
\end{equation}
a function of $n$ and $m$ alone. It does not involve $d$, $e$, or the
modulus $q=m\lcm(d^2,e)$ built from them.

Steps~3 and~4 are therefore unchanged. The factorisation $u=mk$, the
degeneracy lemma, and the unfolding of $\mu^2(k)$ and
$\mathbf 1_{(k,m)=1}$ all act on $k$ through its divisors and not
through its size, and $w_m$ rides along as a coefficient. Two
bookkeeping items move, each by one power of $\log N$: in
Lemma~\ref{lem:degen} a term is now $\ll(\log N)^2$ rather than
$\ll\log N$, which leaves that lemma's $O(N^{o(1)})$ where it was, and
in the two truncation tails of Step~4 the factor $\log N$ standing for
$\max\Lambda$ becomes $(\log N)^2$, so those tails read
$\ll N(\log N)^{3}e^{-\sqrt{\log N}}+N^{1-\theta'/2}(\log N)^2$ and
\[
 \ll N(\log N)^4\log\log N\,e^{-\sqrt{\log N}}
      +M e^{\sqrt{\log N}}(\log N)^3,
\]
still $\ll Ne^{-c\sqrt{\log N}}$ for every fixed $c<1$. The moduli are the same
moduli, so \eqref{eq:level} is untouched.

Step~5 is where the factor costs something, and it costs exactly one
$\log N$.  If $\Tcut_m(t)<1$, the original inner sum is empty; such an
endpoint case can occur for at most one integer $m$, and its omitted
main-term contribution is $O(N^{o(1)})$.  We therefore restrict here to
$\Tcut_m(t)\ge1$. Fix $m$ and write
$\psi_w(y;q,a)=\sum_{n\le y,\ n\equiv a\,(q)}\Lambda(n)w_m(n)$. On
$1\le n\le\Tcut_m(t)$ one has $N-n\ge mK$, so $w_m\ge\log K>0$, and
$w_m$ decreases with $-w_m'(y)=1/(N-y)\ge0$; Abel summation gives
\[
  \psi_w(\Tcut_m;q,N)
   \;=\;w_m(\Tcut_m)\,\psi(\Tcut_m;q,N)
     \;+\;\int_{1}^{\Tcut_m}\frac{\psi(y;q,N)}{N-y}\,dy .
\]
Replacing $\psi(y;q,N)$ by $y/\varphi(q)$ in both places costs at most
\[
  \Bigl(w_m(\Tcut_m)+\int_{1}^{\Tcut_m}\frac{dy}{N-y}\Bigr)
  \max_{y\le N}\Bigl|\psi(y;q,N)-\frac{y}{\varphi(q)}\Bigr|
  \;=\;w_m(1)\,\mathcal E_q\;\le\;(\log N)\,\mathcal E_q,
\]
the bracket telescoping to $w_m(1)=\log((N-1)/m)$. Summing over the
triples $(m,d,e)$ exactly as in Step~5 and applying
Lemma~\ref{lem:BV} with $B=3$ and with $A+1$ in place of $A$ turns
that into $O(N(\log N)^{-A})$. The one extra $\log N$ is absorbed by
the arbitrariness of $A$, which is why no exponent moves.

What is left in place of $\Tcut_m/\varphi(q)$ is
$\varphi(q)^{-1}\bigl(w_m(\Tcut_m)\Tcut_m+\int_1^{\Tcut_m}y(N-y)^{-1}dy
\bigr)=\varphi(q)^{-1}\bigl(\int_1^{\Tcut_m}w_m(y)\,dy+w_m(1)\bigr)$,
again by parts; the $w_m(1)$ is $\ll\log N$ per class and, over the
$\ll MD_0E_0N^{o(1)}$ triples, is $\ll N^{1-\theta'+o(1)}$. Step~6 is
unchanged, the density $c(m)$ of Lemma~\ref{lem:density} being a
property of the modulus and not of the weight.  Here the weighted
integral is $\ll N\log N$, so the two truncation errors in
\eqref{eq:ctrunc} contribute together
\[
  \ll N\log N\left(
       \frac{1}{D_0}\sum_{m<M}\frac{d(m)}{\varphi(m)}
       +\frac{\log\log M}{E_0}\sum_{m<M}\frac{d(m)}{\varphi(m)}
       \right)
  \ll N(\log N)^3(1+\log\log M)e^{-\sqrt{\log N}},
\]
which is $\ll Ne^{-c\sqrt{\log N}}$ for every fixed $c<1$.
Substituting $v=N-y$ in
$\int_1^{\Tcut_m}w_m(y)\,dy$ and taking the actual endpoint
$t=N-1$, for which $\Tcut_m=N-mK$ (since $mK\ge K>1$ for large
$N$), the outcome is the main term
\[
  \AAA(N)\sum_{m<M}\frac{\mu(m)\lambda(m)}{m}
  \int_{mK}^{N}\log(v/m)\,dv.
\]
Substituting $v=mt$ evaluates the integral as
$N\log(N/m)-N-mK\log K+mK$.  The cancellation estimates must be
used here with the same small-$x$ qualification as in
Proposition~\ref{prop:MT}.  Since $M\asymp N^{1-\theta'}$ is a fixed
power of $N$, Lemma~\ref{lem:mu} gives
$G(M)\ll e^{-c\sqrt{\log N}}$, so the $N\log N$ and $-N$ pieces are
$O(Ne^{-c'\sqrt{\log N}})$.  For
$S(M):=\sum_{m\le M}f(m)$, partial summation gives
\[
  S(M)=M G(M)-\int_1^M G(x)\,dx .
\]
Split the integral at $x_0=e^{4C\sqrt{\log N}}$ as in the proof of
Proposition~\ref{prop:MT}: use the trivial $G(x)\ll\log x$ below
$x_0$ and Lemma~\ref{lem:mu} above it.  The same estimates used there
give $K S(M)\ll Ne^{-c'\sqrt{\log N}}$, which disposes of the
$-mK\log K$ and $mK$ pieces.  Finally, with
$\widetilde G(M):=\sum_{m\le M}f(m)\log m/m$ and
$\widetilde G(1)$ denoting the value at $s=1$ of the differentiated
Dirichlet series below, summation by parts on $[M,\infty)$ gives
\[
 |\widetilde G(1)-\widetilde G(M)|
 \le |G(M)|\log M+\int_M^\infty\frac{|G(x)|}{x}\,dx
 \ll e^{-c'\sqrt{\log N}},
\]
because both terms of Lemma~\ref{lem:mu} are integrable on this tail.
Thus only the $-N\log m$ piece contributes to the limiting constant,
and what remains is $-\AAA(N)\,N\,\widetilde G(1)$ with
$\widetilde G(1)=\lim_x\sum_{m\le x}\mu(m)\lambda(m)
\mathbf 1_{(m,N)=1}\log m/m$. The constant is fixed by
\begin{equation}\label{eq:Gtilde}
  \AAA(N)\,\widetilde G(1)\;=\;-\SS(N),
\end{equation}
so that the residual main term is $+\SS(N)\,N+O(N(\log N)^{-A})$.

The sign in \eqref{eq:Gtilde} is forced, and it is the whole identity:
with $F(s)=\sum_m f(m)m^{-s}=\zeta(s)^{-1}H(s)$ as in
Lemma~\ref{lem:mu} one has $\widetilde G(1)=-F'(1)=-H(1)<0$, while
$\AAA(N)H(1)=\SS(N)$ identically --- the local factors pair as
\[
  \Bigl(1-\tfrac{1}{p(p-1)}\Bigr)
  \Bigl(1-\tfrac{1}{(p^2-p-1)(p-1)}\Bigr)
  =1-\tfrac{1}{(p-1)^2}
\]
at $p\nmid N$, and as $p/(p-1)$ at $p\mid N$.

\medskip\noindent
\emph{(ii) The subtracted mean term.} It carries the factor
\[
  B_{\log}(K):=\sum_{\substack{k<K\\ (k,N)=1}}
     \frac{\mu(k)\log k}{\varphi(k)}
  \;=\;-\SS(N)+O\bigl(e^{-c\sqrt{\log K}}\log K\bigr),
\]
which follows from Huang--Li's Lemma~1: subtracting
$\sum_{k\le K,\,(k,N)=1}\mu(k)\varphi(k)^{-1}\log(K/k)
 =\SS(N)+O(e^{-c_1\sqrt{\log K}})$ from
$\log K\cdot\sum_{k\le K,\,(k,N)=1}\mu(k)/\varphi(k)
 =O(\log K\,e^{-c_1\sqrt{\log K}})$ gives the claim.  If $K$
is an integer, replacing $k\le K$ by the defining range $k<K$ changes
this by one term of size $K^{-1+o(1)}\log K$, which is absorbed by the
same error. The coprimality
condition is carried throughout: without it the limiting constant is
not $\SS(N)$. Independently:
with $f(s)=\prod_{p\nmid N}\bigl(1-\frac{p^{-s}}{p-1}\bigr)$ one has
$f(s)=\zeta(s+1)^{-1}h(s)$ with $h(0)=\SS(N)$, so
$B_{\log}(\infty)=-f'(0)=-\SS(N)$. Since the mean term is
$-C(N)B_{\log}(K)$ with $C(N)=\sum_{n<N}\Lambda(n)\mu(N-n)$, it
contributes $+\SS(N)\,C(N)$.

\medskip
Collecting \eqref{eq:goldbachback}, (i) and (ii):
\[
  E_3(\alpha)=\sum_{n<N}\Lambda(n)\Lambda(N-n)
    -\SS(N)N+\SS(N)C(N)+O_A\!\left(\frac{N}{(\log N)^{A}}\right),
\]
which proves Theorem~\ref{thm:C}.  Rearranging, the bound
$E_3(\alpha)\ll_A N(\log N)^{-A}$ is equivalent, at this error scale, to
\[
  \tilde r(N)=\SS(N)\bigl(N-C(N)\bigr)+O_A\!\left(N(\log N)^{-A}\right),
\]
which is Huang--Li's equation~(22).  Their lower bound
\eqref{eq:HLthm1} is a subsequent consequence of that equation.  For the positivity step one needs the quantitative gap, not merely
$\AAA(N)<1$.  The Bombieri--Vinogradov estimate used by Huang--Li and
recalled in the proof of Proposition~\ref{prop:onesided} gives
$|C(N)|\le U(N)=\AAA(N)N+O_B(N(\log N)^{-B})$, while the least prime
$q_0\nmid N$ satisfies $q_0\ll\log N$ and hence
$1-\AAA(N)\gg(\log N)^{-2}$.  Choosing the error exponents in (22) and in $U(N)$ larger than $2$
therefore gives the quantitative lower bound
\[
 \widetilde r(N)\gg \frac{N}{(\log N)^2}.
\]
On the other hand,
\[
 \widetilde r(N)=
 \sum_{\substack{p+q=N\\ p,q\ {\rm prime}}}\log p\,\log q
 +O\bigl(N^{1/2}\log^2N\bigr),
\]
because the omitted terms have at least one proper prime power.  The
error is $o(N(\log N)^{-2})$, so the prime--prime weighted sum is
positive for all sufficiently large even $N$.  This is the binary
Goldbach implication used here.  Moreover, since
$\SS(N)\ge2\prod_{p>2}(1-(p-1)^{-2})\gg1$, equation~(22) gives
$\tilde r(N)\sim\SS(N)N$ if and only if $C(N)=o(N)$; that latter
condition is not proved here. \qed


% =====================================================================
\section{Two weaker sufficient conditions}\label{sec:strength}

\begin{proof}[Proof of Proposition~\ref{prop:onesided}]
From Theorem~\ref{thm:C},
\[
 \widetilde r(N)=\SS(N)(N-C(N))+E_3(\alpha)
 +O_A\!\left(N(\log N)^{-A}\right).
\]
By the triangle inequality,
$|C(N)|\le U(N):=\sum_{n<N}\Lambda(n)\mu^2(N-n)$, and the
Bombieri--Vinogradov calculation in \cite[\S3.3]{HL} gives, for every
$B>0$,
\[
 U(N)=\AAA(N)N+O_B\!\left(N(\log N)^{-B}\right).
\]
Therefore
\[
 \widetilde r(N)\ge \SS(N)(1-\AAA(N))N+E_3(\alpha)
 +O_B\!\left(\SS(N)N(\log N)^{-B}\right)
 +O_A\!\left(N(\log N)^{-A}\right).
\]
The least prime $q_0\nmid N$ satisfies $q_0\ll\log N$, hence
\[
 1-\AAA(N)\ge \frac1{q_0(q_0-1)}\gg(\log N)^{-2},
 \qquad \SS(N)\ge 2C_2>0.
\]
Substituting \eqref{eq:onesided} and taking $A,B>2$ leaves a positive
main term of size at least a fixed multiple of $N(\log N)^{-2}$ for
all sufficiently large $N$.  Since the contribution to
$\widetilde r(N)$ from terms with a proper prime power is
$O(N^{1/2}\log^2N)=o(N(\log N)^{-2})$, the prime--prime weighted
Goldbach sum is positive.  Hence binary Goldbach follows.
\end{proof}

\begin{proof}[Proof of Proposition~\ref{prop:nolog}]
Since $|E_3(\alpha)|\le B_\mu(N)$, the same calculation gives
\[
 \widetilde r(N)\ge \SS(N)(1-\AAA(N))N-B_\mu(N)
 +o\!\left(\SS(N)(1-\AAA(N))N\right).
\]
Equation~\eqref{eq:nolog} leaves a positive margin
$\gg_\varepsilon N(\log N)^{-2}$.  As above, this dominates the
$O(N^{1/2}\log^2N)$ proper-prime-power contribution, so the
prime--prime weighted Goldbach sum is positive for all sufficiently
large even $N$.

Finally, the full hypothesis $EH_\mu(N^{\theta'})$ gives, for every
$L>0$,
\[
 B_\mu(N)\le (\log K)
 \sum_{\substack{k<K\\(k,N)=1}}|\Emu(N;k)|
 \ll_L \frac{N}{(\log N)^{L-1}}.
\]
Taking $L>3$ and using
$\SS(N)(1-\AAA(N))N\gg N(\log N)^{-2}$ shows that
\eqref{eq:nolog} holds eventually, for example with any fixed
$0<\varepsilon<1$ after increasing $N$.
\end{proof}

% =====================================================================
\section{The correction \texorpdfstring{$\Delta$}{Delta} to
equation~(18)}\label{sec:delta}

The arXiv versions v1 and v2 of \cite{HL} replace an
$n$-dependent cutoff in the derivation of their equation~(18) by the
fixed outer range $k<(N-1)/\alpha$.  We do not claim discovery of this
defect: it had been reported to the authors earlier, and Moya
\cite{Moya} subsequently gave a finite counterexample and a repair that
keeps the moving endpoint.  The purpose here is narrower.  We write the
difference between the fixed-cut and moving-cut formulations as an
explicit term $\Delta$, bound it, and use the same estimate to transfer
Theorem~\ref{thm:A} to the corrected cutoff.

Huang--Li define
\[
  S_2(\alpha)=\sum_{n<N}\Lambda(n)\,\mu^2(N-n)\,
     \tilde\Lambda_\alpha(N-n),
  \qquad
  \tilde\Lambda_\alpha(u)=\sum_{d\mid u,\ d>\alpha}\mu(d)\log(1/d),
\]
and substitute $k=u/d$, so that the constraint $d>\alpha$ becomes
\begin{equation}\label{eq:constraint}
  k\;<\;\frac{N-n}{\alpha},
\end{equation}
an \emph{$n$-dependent} bound. Their displayed equation~(18) reads
\[
  S_2(\alpha)\;=\;\sum_{k<\frac{N-1}{\alpha}}\mu(k)
    \sum_{\substack{n<N\\ n\equiv N\ (k)}}
    \Lambda(n)\mu(N-n)\log\Bigl(\frac{k}{N-n}\Bigr),
\]
in which \eqref{eq:constraint} has been replaced by the $n$-free bound
$k<(N-1)/\alpha$. The two differ. Writing $N-n=mk$, the condition
\eqref{eq:constraint} is $m>\alpha$, so the right-hand side of~(18)
contains, in addition to $S_2(\alpha)$, exactly the terms with
$m\le\alpha$. On those terms
$\mu(k)\mu(N-n)=\mu(m)\mu^2(k)\mathbf 1_{(k,m)=1}$ and
$\log(k/(N-n))=-\log m$, so
\begin{equation}\label{eq:Delta}
  S_2(\alpha)\;=\;\text{RHS of (18)}\;+\;\Delta,
  \qquad
  \Delta\;=\;\sum_{2\le m\le\alpha}\mu(m)\log m
    \sum_{\substack{k<(N-1)/\alpha\\ (k,m)=1}}\mu^2(k)\,\Lambda(N-mk).
\end{equation}
(The term $m=1$ drops out since $\log1=0$.) The trivial bound is
$\Delta\ll N(\log N)^2$, which is far above the target
$O(N(\log N)^{-A})$.  Thus the trivial estimate does not make
$\Delta$ negligible, and the term needs its own treatment.

It does close, by the machinery already used above and under
hypotheses Huang--Li already assume. Indeed $\Delta$ has exactly the
shape of the residual \eqref{eq:R1}: the M\"obius factor sits on the
\emph{short} variable $m\le\alpha$, and the long variable $k$ carries
only $\mu^2\ge0$. Since $\mu(m)\neq0$ on the range of summation, the
two conditions on $k$ collapse into one,
\begin{equation}\label{eq:mu2mk}
  \mu^2(k)\,\mathbf 1_{(k,m)=1}\;=\;\mu^2(mk),
\end{equation}
and expanding $\mu^2(mk)=\sum_{b^2\mid mk}\mu(b)$ with the truncation
$b\le B_1:=(\log N)^{A+4}$ --- the same one Huang--Li use in their
\S3.1, where it is called $B$; the subscript keeps it apart from the
free exponent of Lemma~\ref{lem:BV} --- reduces $\Delta$ to sums of
$\Lambda$ over progressions to moduli
$[m,b^2]\le\alpha B_1^2=\alpha(\log N)^{2A+8}$. The tail $b>B_1$ is
discarded as the $d$-tail of Step~4 was, but counted on the pair
$(m,k)$ and not on the modulus alone. Such terms have $b^2\mid mk$,
which holds exactly when $k\equiv0$ to the modulus $b^2/(b^2,m)$; so
for each $m$ the number of admissible $k<K:=(N-1)/\alpha$ is at most
$K(b^2,m)/b^{2}$, and
\[
  \sum_{m\le\alpha}(b^2,m)
   \;=\;\sum_{g\mid b^2}g\cdot\#\{m\le\alpha:(b^2,m)=g\}
   \;\le\;\sum_{g\mid b^2}g\cdot\frac{\alpha}{g}
   \;=\;\alpha\,d(b^2).
\]
Hence, with $\Lambda\ll\log N$, the factor $\log m\le\log\alpha\ll\log
N$ that $\Delta$ carries on the short variable, and $K\alpha=N-1$,
\[
  \ll(\log N)^{2}\!\!\sum_{b>B_1}\frac{K\alpha\,d(b^2)}{b^{2}}
  \;\ll\;N(\log N)^{2}\!\!\sum_{b>B_1}\frac{d(b^2)}{b^{2}}
  \;\ll\;\frac{N(\log N)^{2}(\log B_1)^{2}}{B_1}
  \;\ll\;\frac{N(\log\log N)^{2}}{(\log N)^{A+2}},
\]
which is $\ll_AN(\log N)^{-A}$. The last step is
$\sum_{b>X}d(b^2)b^{-2}\ll(\log X)^2/X$, by partial summation from
$\sum_{b\le Y}d(b^2)\ll Y(\log Y)^2$.

Counting on the pair removes any extra distributional input from this
tail estimate. Bounding each
progression separately by $N/[m,b^2]+1$ with $[m,b^2]\ge b^2$ drops the
factor $1/m$ and drops the constraint that an admissible $k$ exists at
all; with the factors $\Lambda\ll\log N$ and
$\log m\ll\log N$, the resulting $+1$ then contributes
$\ll\alpha\sqrt N(\log N)^2$.  That crude estimate would therefore
require $\alpha<\sqrt N$ up to logarithmic factors; the
pair count above avoids this auxiliary restriction. What is left is the progression main term.  For fixed $m$ the $k$-range
has length $K=(N-1)/\alpha$, hence the corresponding $n=N-mk$ interval
has length $mK$.  Multiplying this by the local density
$\AAA(N)\lambda(m)/m$ gives
\[
 \AAA(N)K\sum_{m\le\alpha}\mu(m)\lambda(m)\log m,
\]
up to the already bounded progression errors.  Partial summation from
Lemma~\ref{lem:mu} makes this $O(Ne^{-c\sqrt{\log N}})$.

The density here is not the one Lemma~\ref{lem:density} proves. That
lemma evaluates the first of
\[
  \sum_{d\ge1}\sum_{e\mid m}
    \frac{\mu(d)\mu(e)\mathbf 1_{(d,N)=1}}{\varphi(m\lcm(d^2,e))},
  \qquad
  \sum_{b\ge1}\frac{\mu(b)}{\varphi([m,b^2])},
\]
and what the expansion above leaves is the second, over the moduli
restricted to $(\,\cdot\,,N)=1$. The two agree, and the check is the same local
computation: for squarefree $m$ the factor of the second is
$\tfrac{1}{p-1}-\tfrac{1}{p(p-1)}=\tfrac1p$ at $p\mid m$ and
$1-\tfrac{1}{p(p-1)}$ at $p\nmid m$, and the restriction sets the
latter to $1$ at $p\mid N$ exactly as in Lemma~\ref{lem:density} ---
the same coprimality point used above. The
product is $\AAA(N)\lambda(m)/m$, as written.

The resulting moduli stay within the level supplied by the
corresponding Huang--Li hypothesis. Hence:
\begin{itemize}
\item In the Corollary-1 regime, $\alpha=N^{1-\theta'}<N^{1/2}$ and
  Bombieri--Vinogradov bounds $\Delta$ without an additional
  Elliott--Halberstam hypothesis.
\item In the general Theorem-1 regime, $\alpha\asymp N^{\theta}$ and
  $\Delta$ closes under the hypothesis
  $EH(N^{\theta}(\log N)^{2A+8})$, which is assumed there.
\end{itemize}
Thus the cutoff defect does not by itself invalidate the stated
conditional conclusions: Moya repairs the moving-cut argument under the
original hypotheses, while the outside formulation above controls the
extra term in the same regimes.  In the corrected moving-cut formulation
$\Delta$ does not arise.

\subsection{The two formulations differ by \texorpdfstring{$\Delta$}%
{Delta}}\label{sec:movingcut}

The corrected argument keeps the moving cut: with
$Y_k=\lceil N-\alpha k\rceil-1$ the inner sum stops at $Y_k$, the terms
with $m\le\alpha$ never enter, and \eqref{eq:Delta} has nothing to
correct. Every statement above was proved against the fixed cut, so one must
check separately what survives the move.  In the Corollary-1 regime the
following proposition gives that transfer, and the term carrying the
difference has the same short-variable shape as $\Delta$.

Write $J=T_1(N-1)$ for the fixed-cut object at its extreme endpoint and
$T_1^\ast=\sum_{k<K,(k,N)=1}\mu(k)\Emu(Y_k;k)$ for the moving-cut one.
The difference is carried by the $n$ the moving cut drops, and there
\[
  n>Y_k \iff N-n\le\alpha k \iff m\le\alpha,
  \qquad m=\frac{N-n}{k},
\]
so on exactly those terms $\mu(k)\mu(N-n)=\mu(k)\mu(mk)$, which
vanishes unless $mk$ is squarefree and equals $\mu(m)\mu^2(k)$
otherwise: the M\"obius factor on the \emph{short} variable, the long
one carrying only $\mu^2\ge0$. That is \eqref{eq:Delta}'s shape with
weight $1$ in place of $\log m$.

\begin{proposition}[the bound survives the corrected cut]%
\label{prop:movingcut}
Fix $\theta'\in(1/2,1)$, put $\alpha=N^{1-\theta'}$ and
$K=(N-1)/\alpha$. Then for every $A>0$,
$T_1^\ast\ll_{A,\theta'}N(\log N)^{-A}$, unconditionally.
\end{proposition}

\begin{proof}
Exchanging as above,
$J-T_1^\ast=\Xi_{\mathrm{main}}-\Xi_{\mathrm{mean}}$ with
\[
  \Xi_{\mathrm{main}}=\sum_{m\le\alpha}\mu(m)\!\!
    \sum_{\substack{k<K\\ (k,mN)=1}}\!\!\mu^2(k)\,\Lambda(N-mk),
  \qquad
  \Xi_{\mathrm{mean}}=\!\!\sum_{\substack{k<K\\(k,N)=1}}\!\!
    \frac{\mu(k)}{\varphi(k)}\bigl(C(N-1)-C(Y_k)\bigr),
\]
the constraint $n\ge1$ being automatic since $m\le\alpha$ and $k<K$
force $mk<N-1$.

$J\ll_{A,\theta'}N(\log N)^{-A}$ is Theorem~\ref{thm:A} at $t=N-1$,
which lies in the range of its supremum, and it is used exactly where
Theorem~\ref{thm:A} needs it. It is one of the two places this proof
needs $\theta'>1/2$; the other is the bridge term below, whose moduli
carry the same inequality.

In $\Xi_{\mathrm{main}}$ the term $m=1$ is
$\sum_{k<K,(k,N)=1}\mu^2(k)\Lambda(N-k)\le\sum_{N-K<j<N}\Lambda(j)\ll
K\log N=N^{\theta'}\log N$, which is $\ll_AN(\log N)^{-A}$ because
$\theta'<1$. The terms $2\le m\le\alpha$ are \eqref{eq:Delta} with
$\log m$ replaced by $1$ and close by the argument above, the main term
dying by Lemma~\ref{lem:mu} applied directly --- the weight being $1$,
no partial summation is needed to remove a $\log m$. One step is not
verbatim: there the inner variable carries $(k,m)=1$ and
\eqref{eq:mu2mk} absorbs it into $\mu^2(mk)$, while here it carries
$(k,mN)=1$, and the $(k,N)=1$ that is left over after \eqref{eq:mu2mk}
is discarded by Lemma~\ref{lem:degen} at a cost of $O(N^{o(1)})$, as at
\eqref{eq:R1}.

For $\Xi_{\mathrm{mean}}$, put $K_-:=\lceil K\rceil-1$, so that
for integer $k$ the condition $k<K$ is exactly $k\le K_-$. Since
$n>Y_k$ iff $k>R_n:=\lceil(N-n)/\alpha\rceil-1$, exchanging gives
$\Xi_{\mathrm{mean}}=\sum_{n<N}a_n(\rho_N(K_-)-\rho_N(R_n))$, where
\[
  \rho_N(R):=\sum_{\substack{d\le R\\ (d,N)=1}}\frac{\mu(d)}{\varphi(d)}
  \;\ll\;e^{-c_1\sqrt{\log R}}
  \qquad(\log N\ll\log R)
\]
by \cite[Lemma~1]{HL}, a form of \cite[Lemma~2.1]{GY} --- the
hypothesis $\log N\ll\log R$ is part of the statement and is checked
below --- and $a_n=\Lambda(n)\mu(N-n)$. Split at $N-n=H$ with
$H:=2N^{1-\theta'/2}$. On $N-n\le H$ use the elementary bound
$|\rho_N(R)|\le\sum_{d\le R}1/\varphi(d)\ll\log(2+R)\ll\log N$
and $\sum_{N-H\le n<N}|a_n|\ll H\log N$.  This contributes
$\ll H(\log N)^2\ll_AN(\log N)^{-A}$ since
$1-\theta'/2<1$. On $N-n>H$ one has $R_n\ge N^{\theta'/2}$, so
both arguments $K_-\asymp K$ and $R_n$ of $\rho_N$ meet its hypothesis
against $N$, and both
terms are $\ll e^{-c\sqrt{\theta'\log N}}$; with
$\sum_{n<N}|a_n|\ll N$ this is $\ll_AN(\log N)^{-A}$ as well.
\end{proof}


% =====================================================================
\section{Relation to prior work}\label{sec:lit}
Murty and Vatwani \cite{MV17} use the same convolutional mechanism in
the shifted-prime setting, and Huang--Li \cite{HL} adapt it to the
Goldbach shift.  Theorem~\ref{thm:C} claims no novelty for that chain;
it keeps the surviving $E_3$ term explicit.  The fixed-class feature is
also present in \cite{MV17,HL}.

More recently, Cantarini \cite{Cantarini} studies weighted averages of
the diagonal M\"obius-twisted Elliott--Halberstam problem, under GRH and
additional hypotheses, for Sobolev and H\"older--Zygmund weights.  That
work concerns weighted diagonal averages over moduli.  The present
Theorem~\ref{thm:A} instead exploits the particular signed fixed-class
combination occurring in Huang--Li's $E_4$ term and, in the
Corollary-1 regime, obtains its required bound unconditionally from
divisor switching and Bombieri--Vinogradov.  We make no claim that this
mechanism supplies a corresponding bound for the full diagonal
$EH_\mu$ average studied there.

The moving-cut defect is not claimed as new: Moya \cite{Moya} records
it and repairs the original argument under Huang--Li's hypotheses.
Section~\ref{sec:delta} gives a complementary outside correction.  The
net result is therefore limited: the $E_4$ use of $EH_\mu$ is removed
in the regime studied, while the surviving $E_3$ correlation remains
unestimated.  No unconditional Goldbach conclusion is claimed.

% =====================================================================
% =====================================================================
\begin{thebibliography}{99}\bibitem{GY} D.~Goldston and C.~Y{\i}ld{\i}r{\i}m,
\emph{Higher correlations of divisor sums related to primes~I},
Integers \textbf{3} (2003), A5.

\bibitem{HL} J.-J.~Huang and H.~Li, \emph{On the connection between
the Goldbach conjecture and the Elliott--Halberstam conjecture},
arXiv:2005.03811 [math.NT]; v1 (May 2020), v2 (August 2022). Also published in a Springer volume,
\texttt{doi:10.1007/978-3-030-67996-5\_17}; references here are to the
arXiv version.

\bibitem{MV17} M.~R.~Murty and A.~Vatwani, \emph{Twin primes and
the parity problem}, J. Number Theory \textbf{180} (2017), 643--659.

\bibitem{MoV} H.~L.~Montgomery and R.~C.~Vaughan,
\emph{Multiplicative number theory~I: classical theory}, Cambridge
Studies in Advanced Mathematics \textbf{97}, Cambridge University
Press, 2007.

\bibitem{Cantarini} M.~Cantarini,
\emph{Averages of diagonal Elliott--Halberstam problem twisted by
M\"obius function with Sobolev and H\"older--Zygmund weights},
arXiv:2607.09110 [math.NT], July 2026.

\bibitem{MoV2} H.~L.~Montgomery and R.~C.~Vaughan,
\emph{Multiplicative number theory~II: primes and sieves}, Cambridge
Studies in Advanced Mathematics \textbf{218}, Cambridge University
Press, 2026.  The maximal Bombieri--Vinogradov theorem used here is
from \S20.2.

\bibitem{Moya} R.~Moya, \emph{A Moving-Cut Correction to the
Huang--Li Conditional Goldbach Argument: Consequences for Diagonal
M\"obius-Twisted Elliott--Halberstam Hypotheses}, Zenodo, August 2026;
DOI \texttt{10.5281/zenodo.21939906}. Gives a finite counterexample to the fixed-cut rearrangement and a
moving-endpoint repair under Huang--Li's original hypothesis.

\end{thebibliography}

% =====================================================================
\paragraph{Disclosure statement.} No potential competing interest, financial or non-financial, was
reported by the author. No funding was received for this work. AI-assisted
drafting and computational auditing were used in preparing the manuscript;
no computation or AI-generated output is used as a premise of a proof.

\paragraph{Data availability statement.} The audit code, stored outputs, and Lean~4 development accompanying this
manuscript are openly available at
\url{https://github.com/soke91/arithmetic-convolution-fields}, in the
directory \texttt{P1-mobius-fixed-class/}. No data derived from human or
animal subjects were used.

% =====================================================================
\end{document}
